%% ============================================================
%% VALIDATION PROTOCOL
%% Gate values in the prose must match tab_val_targets. The return-temperature gate is
%% 2.5 K -- the literature value, fixed before calibration. If the table shows 1.0 K,
%% the table is stale.
%% ============================================================

The network was upgraded with the heat pump, electrode boiler and thermal storage
\emph{after} the measurement record was collected, so validation follows a two-stage
strategy after \citet{Kus2025} and \citet{Maldonado2024}.

\paragraph{Stage 1: direct network validation.}
The pipe model is validated against pre-upgrade monitoring data with the new assets set to
zero capacity, calibrated by the branch-sequential $U_p$-tuning procedure of
\citet{Maldonado2024} but constrained to a defensible coefficient range for the pipe class.
Table~\ref{tab:valtargets} lists the key performance indicators and their gates. The gates
are taken from the values reported for comparable validations by \citet{Maldonado2024} and
\citet{Kus2025} and were fixed before the calibration was run, so they are not thresholds
chosen to accommodate the result: the return-temperature gate is the $2.5$\,K mean absolute
error of that literature, not a tolerance widened after the fact.

% CHANGED: "the temperature gates" -> the supply-temperature gates specifically.
% The return-temperature gate is measured at the source, where instrumentation is
% adequate; the nodal supply gates are the metering-limited ones.
How those gates are used is deliberate: they are reported in full in
Section~\ref{subsec:validation}, including the $\langle$N$\rangle$ that are not met. The
nodal supply-temperature gates in particular cannot be met by any model validated against
this instrumentation, for the reason set out there: the consumer sensors are billing meters
downstream of mixing valves and do not read the primary junction temperature. Retaining
gates the instrumentation cannot support, and reporting the failures against them, is the
point rather than an oversight: it is the evidence for the claim that a model resolved more
finely than its metering cannot have that resolution verified. The conclusions are anchored
instead to the annual delivered energy, and hence the annual network loss, which is gated,
met, and the quantity the cost results depend on.

\paragraph{Stage 2: necessary-condition verification.}
Dispatch of the new assets cannot be validated against measurements, so it is checked
against plausibility bounds. Two of these are enforced constraints and their satisfaction
confirms only that the implementation matches the formulation: the storage state of charge
within $[0.05, 0.95]\bar{E}_s$ and the exclusion of simultaneous charge and discharge. Two
are unconstrained outputs and therefore carry information: the realised heat-pump
coefficient of performance, which stays within $[2.5, 5.5]$ across all hours without being
bounded there, and per-step energy closure, which holds to machine precision. These are
necessary conditions, not sufficient ones; a post-upgrade measurement campaign is the
natural follow-up and is listed among the limitations.

\paragraph{Structural sanity checks.}
Two checks confirm internal consistency. First, the nonlinear reference with its quadratic
constraints deactivated reproduces the baseline objective to within the $0.01$\,\% solver
tolerance, so the two formulations are the same programme up to the gap rather than merely
close. Second, $\mathrm{cost}(\CP) \le \mathrm{cost}(\Lone)$ holds on all 135 synthetic
instances and on Memmingen. This ordering is expected on structural grounds: relative to
$\Lone$, the copperplate both enlarges the feasible set, by dropping the nodal transport
constraints, and removes a non-negative term from the objective, the network loss that must
otherwise be generated. Both act in the same direction, so the inequality is a property of
the formulation and a violation would indicate an implementation error rather than a finding.